\section{Gravity}

Gravity is not a force added to the mesh. It is the \emph{shape of the mesh} acting back on what moves through it. A body bends the weave of docks around it, and other bodies, sliding through that bent weave, fall toward the body. There is nothing to add. The geometry already there is the attraction.

The vibe mesh reaches gravity on \emph{two faces}, and they must be kept apart. The flat three-dimensional cusp gives a clean Newtonian $1/r$ potential, measured. The curved hyperbolic bulk should give curvature-sourced attraction, and that part is an \emph{honest open}. We state both, and we do not blur them.

\begin{principle}[The two faces, kept apart]
The flat cusp gives a clean $1/r$ Newtonian potential. Curved-bulk gravity on the hyperbolic mesh did not discriminate cleanly. These are two different objects. The clean result is the flat-cusp one. We never let the flat-cusp success stand in for the curved-bulk question.
\end{principle}

This section separates what is \emph{measured}, what is \emph{derived from measured inputs}, and what is \emph{open}. We begin by naming the two objects, then build the discrete attraction field that completes a self, then turn to the clean flat-cusp Newtonian law, the holographic correlator, the route to the full Einstein equations, gravitational waves, and finally the honest frontier that remains.

\subsection{The two objects, kept apart}

Before any derivation it pays to lay the two gravitational objects side by side, with their geometry, their falloff, and their status, so they are never confused later.

\begin{table}[ht]
\centering
\begin{tabular}{@{}llll@{}}
\toprule
object & where & falloff & status \\
\midrule
Newtonian potential & the flat 3D cusp & $1/r$ & measured (difference method) \\
bulk-mediated correlator & through the curved bulk & $1/r^2$ & derived plus measured \\
curved-bulk attraction & the hyperbolic bulk & no clean law & open \\
\bottomrule
\end{tabular}
\caption{The three gravitational objects of the mesh. The clean $1/r$ Newton result is a flat-cusp result. The curved-bulk pull is a separate question, and it is the one that did not resolve.}
\label{tab:gravity-objects}
\end{table}

The clean $1/r$ Newton result lives in the \emph{flat 3D cusp}. The curved-bulk pull on the hyperbolic mesh is a separate question, and it is the one that did not resolve. Do not conflate them.

\subsection{The discrete gravity field}

The base is \emph{discrete and deterministic}. Eight atoms, ternary tone, integer beats. The gravity field must be discrete and deterministic too. No real numbers, no decimals. A real-valued potential would smuggle the continuum back in through a side door, and the whole point of the mesh is that it does not.

\begin{definition}[The discrete gravity field]
A per-dock \emph{integer} well, dug by mass, that bodies slide down by one site step per beat. The well is a bounded small integer, the count of at most $\mathrm{cap}$ token-bits in a dock. Mass sources the well, the well guides the motion, and the motion is the attraction.
\end{definition}

The field has four pieces. We name them, then state the relation that ties them together.

\begin{table}[ht]
\centering
\begin{tabular}{@{}lll@{}}
\toprule
piece & what it is & bound \\
\midrule
mass & a rest charge, one per dock at most & $0$ or $1$ occupancy \\
source & occupied dock with $3+$ occupied neighbours (bulk only) & $0$ or $1$ \\
potential & per-dock integer well & clamped to $[-\mathrm{cap}, \mathrm{cap}]$ \\
force & hop to the lowest-well empty neighbour site & one step per beat \\
\bottomrule
\end{tabular}
\caption{The four pieces of the discrete gravity field.}
\label{tab:gravity-pieces}
\end{table}

Mass is a \emph{rest charge}, a tone sitting still rather than streaming along one of the $24$ sites. At most one rest charge per dock, so the mass field is a plain $0$-or-$1$ occupancy. A dock is a \emph{source} only when it is occupied and has at least three occupied neighbours among its $24$ sites. This bulk-only rule is what removes the self-force, taken up below.

The potential is a per-dock integer $\phi$, bounded to $[-\mathrm{cap}, \mathrm{cap}]$. It satisfies a \emph{discrete Poisson relation}, solved by relaxation. One relaxation sweep is, for every dock $c$,
\[
\phi_{\text{next}}[c] \;=\; \mathrm{clamp}\!\left(\, \mathrm{round}\!\left( \frac{1}{24}\sum_{n \in N(c)} \phi[n] \right) - \mathrm{strength}\cdot \mathrm{source}[c],\; -\mathrm{cap},\; +\mathrm{cap} \,\right),
\]
where $N(c)$ is the set of $24$ site neighbours of dock $c$. The neighbour average diffuses the well into a smooth bowl. The source term digs the well deeper at each mass dock. The clamp keeps $\phi$ inside a few bits. Iterating this sweep a few times relaxes the well to the body's current shape.

The clamp is the key minimality point.

\begin{itemize}
\item The well is a \emph{bounded small integer}, a count of at most $\mathrm{cap}$ token-bits in a dock.
\item With $\mathrm{cap} = 6$ there are thirteen levels, about \emph{four bits per dock}.
\item An unbounded integer would be a real number in disguise. The bound forbids that.
\item The well \emph{depth} is capped. The force \emph{range} grows with $\mathrm{cap}$ and with the number of sweeps.
\end{itemize}

How few bits suffice was measured, and it passes. Pure ternary, $\mathrm{cap} = 1$, is too coarse. Its range is about one dock, and it heals only a contact nudge. With $\mathrm{cap} = 6$ the field heals a piece broken off \emph{three docks away}. So the attraction is a small bounded multi-level field, more than ternary, far from an arbitrary integer. The well itself is coarse. Its \emph{gradient}, which neighbour dock is lower, is all the dynamics needs.

\begin{result}[The field is slaved to mass]
The well carries no independent conserved content. It is recomputed each beat from where the mass sits. This is the \emph{Newtonian instantaneous-potential} picture. The field tracks the mass with no propagation delay.
\end{result}

\subsection{No self-force}

Only \emph{bulk} docks are sources. A lone or displaced mass is not a source, so it does not sit in its own well.

\begin{itemize}
\item A displaced piece is a \emph{test particle}. It feels only the body's well, not its own.
\item Without this rule, a displaced mass would sit in the well it digs for itself and never return.
\item This is the discrete analogue of \emph{a particle not self-gravitating}.
\end{itemize}

Bulk masses do not move. They are surrounded, with no empty lower neighbour, and the knit skips them explicitly. Only the surface and the displaced pieces feel a net pull. This is exactly what keeps a body from collapsing inward on itself.

\subsection{Excluded volume is the short-range repulsion}

One mass per dock. The target of any hop must be empty. This single rule is the hard core that balances the attraction.

\begin{itemize}
\item Attraction (down the well) plus excluded volume (a hard core) gives a \emph{stable preferred density}.
\item A body holds its size and repairs a hole, instead of collapsing to a point.
\end{itemize}

\begin{principle}[Attraction plus exclusion gives a body]
Down-well motion pulls mass together. Excluded volume stops two masses from sharing a dock. The balance of the two fixes a preferred density, so the body has a definite size that it defends and repairs.
\end{principle}

\subsection{The per-beat algorithm}

Given the will (the tones) and the carried well $\phi$, one beat runs six steps in order.

\begin{enumerate}
\item \textbf{Collide then stream} the moving charges. This is the base knit. It advances the radiation.
\item \textbf{Source.} Mark each occupied dock with three or more occupied neighbours.
\item \textbf{Relax.} Run a few warm-started sweeps to update the well.
\item \textbf{Move.} Hop each non-bulk mass to the lowest empty neighbour, if it is lower. This step is the attraction.
\item \textbf{Bath.} Absorb at the open wake. Radiation that reached the edge is gone for good.
\item \textbf{Carry} the well to the next beat.
\end{enumerate}

The well is \emph{warm-started}. The body barely moves between beats, so the well is carried over and only a few sweeps run each beat to track the small change. The field is fast relative to the body, the way instantaneous Newtonian gravity is fast relative to the planets it guides.

The parameters and their roles are fixed and small.

\begin{table}[ht]
\centering
\begin{tabular}{@{}lll@{}}
\toprule
name & default & role \\
\midrule
$\mathrm{cap}$ & $6$ (about $4$ bits) & the bound on $\phi$, the well depth, the bit budget \\
$\mathrm{strength}$ & $=\mathrm{cap}$ & the source term, set so the well saturates at the bound \\
sweeps & $\sim 24$ then $4$ & how far the well spreads, the effective range \\
$\mathrm{minNeighbours}$ & $3$ & the bulk threshold, defines the source, kills self-force \\
excluded volume & $1$ mass per dock & the short-range repulsion, sets the body size \\
\bottomrule
\end{tabular}
\caption{The parameters of the discrete gravity field. Each is a small fixed integer with a single clear role.}
\label{tab:gravity-params}
\end{table}

\subsection{What the discrete field does}

The field was measured on a bound body, and the three properties of a self all hold.

\begin{table}[ht]
\centering
\begin{tabular}{@{}ll@{}}
\toprule
property & result \\
\midrule
identity & body persists, no collapse (occupied $89$, extent $4$, unchanged) \\
self-repair & a piece broken off $3$ docks away returns (extent $6$ falls to $4$) \\
radiation & a disturbance sheds to the bath (open difference to $0$), persists on the closed torus (difference $8$) \\
\bottomrule
\end{tabular}
\caption{The measured behaviour of the discrete gravity field on a bound body.}
\label{tab:gravity-measured}
\end{table}

The field is \emph{mass-conserving}. The move step relocates masses, it never creates or destroys them. It is \emph{finite-range}. In four dimensions the well's gradient weakens fast, so it repairs small displacements, not teleported pieces. That is fine for a self, which only ever needs to repair small perturbations.

\subsection{The honest reversibility caveat}

The base is \emph{reversible}. This gravity field, as specified, is \emph{not}.

\begin{itemize}
\item The relaxation sweep is diffusive. It forgets initial conditions.
\item The down-well hop is dissipative. Many configurations map to the same lower-energy one.
\end{itemize}

So this gravity is an \emph{effective instantaneous-Newtonian attraction}. It is a working demonstration that the one missing ingredient is an attraction, and that a discrete attraction suffices. It is not yet a clean reversible base field.

\begin{result}[The remaining gap is reversible discrete gravity]
A reversible discrete gravity would need the well to be a genuine dynamical field with its own conserved degrees of freedom, a propagating wave the mass sources and that back-reacts. The naive second-order integer wave field for this is \emph{unstable}, the same discrete Klein-Gordon instability that killed the sine-Gordon kink. A stable, reversible, discrete, propagating gravity is open. This is the real remaining gap, narrower than the whole self, but real.
\end{result}

\subsection{Clean Newton in the 3D cusp}

In the flat three-dimensional cusp, the static potential of a point source is the \emph{lattice Green's function} of the Laplacian. This is the discrete cousin of the continuum $1/r$ potential, and it is where the clean measured result lives.

The naive finite-patch sum is unreliable. It is contaminated by the $k=0$ zero mode and by boundary effects, both of which corrupt the falloff before it can be read off. The correct tool is the \emph{difference method}.

\begin{principle}[The difference method]
Measure $G(r) - G(r_0)$ with the regular cosine integrand. The subtraction cancels the zero-mode contamination, and the clean $1/r$ falloff emerges.
\end{principle}

The result is clean.

\begin{table}[ht]
\centering
\begin{tabular}{@{}ll@{}}
\toprule
quantity & value \\
\midrule
form & $G(r) = a/r + b$ \\
coefficient $a$ & $\sim 0.083 \approx 1/(4\pi)$ \\
dimension & 3D (the cusp) \\
\bottomrule
\end{tabular}
\caption{The measured flat-cusp Newtonian potential. The coefficient matches the continuum $1/(4\pi)$.}
\label{tab:newton-cusp}
\end{table}

It is $1/r$ \emph{because the cusp is three-dimensional}. The force is dimension-resolved. This is the Newtonian potential falling out of the geometry, not added by hand. This is the flat-cusp result, and it is clean.

\subsection{The bulk-mediated holographic correlator}

Holography is the idea that everything happening in a chunk of gravitating space can be written on its lower-dimensional boundary, the way a flat hologram stores a three-dimensional scene. The mesh realizes this through the curved bulk.

Between two boundary points, the bulk-mediated correlator decays as a clean $1/r^2$. This is a different object from the flat-cusp Newtonian $1/r$, mediated through the curved bulk rather than along the flat cusp.

\begin{table}[ht]
\centering
\begin{tabular}{@{}ll@{}}
\toprule
method & result \\
\midrule
exact Bethe-lattice cavity recursion & $\mu = 1/b$, $\alpha = 2$ \\
validation vs directly-solved finite tree & matches \\
tree-tensor (MERA) contraction & $\sim 1/r^{1.93}$ \\
\bottomrule
\end{tabular}
\caption{The bulk-mediated holographic two-point function, by three methods that agree.}
\label{tab:holo-correlator}
\end{table}

This is the holographic two-point function. It is a \emph{different object} from the flat-cusp Newtonian $1/r$. The Ryu-Takayanagi law and the holographic codes that this correlator sits beside are owned by the holography section.

\subsection{The Einstein equations, from a measured area law}

The mesh reaches \emph{full general relativity} to the level any emergent-gravity program reaches, not only the Newtonian limit. The route is Jacobson's entropic derivation, which turns thermodynamics on local horizons into the gravitational field equation.

\begin{principle}[Einstein from the area law]
Given the measured area law $S = A/4G$ and the Unruh temperature $T = \kappa/2\pi$, the Clausius relation $dQ = T\,dS$ on every local Rindler horizon \emph{is} the full nonlinear Einstein equation.
\end{principle}

That equation is $G_{\mu\nu} + \Lambda g_{\mu\nu} = 8\pi G\, T_{\mu\nu}$. The checks are quantitative, not qualitative.

\begin{table}[ht]
\centering
\begin{tabular}{@{}ll@{}}
\toprule
check & result \\
\midrule
Einstein coefficient & exactly $8\pi$ ($2\pi/\eta = 25.133 = 8\pi$) \\
same $G$ as Newton & yes \\
weak field & Poisson with the correct $4\pi$ \\
light bending & $4GM/b$, the GR factor $2$ (ratio $2.0000$) \\
\bottomrule
\end{tabular}
\caption{Quantitative checks of the emergent Einstein equation against general relativity.}
\label{tab:einstein-checks}
\end{table}

The factor-two light bending is the cleanest observable separating full general relativity from Newton. The metric curves time \emph{and} space, so a photon deflects by twice the naive Newtonian amount. The mesh gets it right, to four figures.

Black-hole thermodynamics is complete and self-consistent.

\begin{table}[ht]
\centering
\begin{tabular}{@{}ll@{}}
\toprule
relation & result \\
\midrule
first law $dM = T\,dS$ & error $\sim 10^{-9}$ \\
Smarr $M = 2TS$ & exact \\
Bekenstein bound & saturated \\
heat capacity & $C = -8\pi M^2$ (evaporation instability) \\
lifetime & $\sim M^3$ \\
\bottomrule
\end{tabular}
\caption{Black-hole thermodynamics on the mesh.}
\label{tab:bh-thermo}
\end{table}

The mesh's own de Sitter horizon carries the same thermodynamics, with $T = H/2\pi$ and $\Lambda = 3H^2$ fixed by the measured growth of the wake. This shared horizon temperature is developed in the cosmology section.

\subsection{Gravitational waves}

A massless \emph{spin-two graviton} lives on the flat cusp. It is not added. It is the transverse-traceless ripple of the same metric the area law measures.

\begin{table}[ht]
\centering
\begin{tabular}{@{}ll@{}}
\toprule
property & value \\
\midrule
dispersion & $\omega = |k|$, $z = 1$ \\
polarizations & exactly $2$ transverse-traceless \\
binary waveform & quadrupole, GW frequency $=$ twice orbital \\
radiated power & $P = (32/5)\,\mu^2 a^4 \omega^6$ \\
inspiral chirp & $f \sim (t_c - t)^{-3/8}$, the LIGO exponent (verified $-0.375$) \\
\bottomrule
\end{tabular}
\caption{The gravitational wave sector on the flat cusp.}
\label{tab:gw}
\end{table}

The inspiral chirp exponent $-3/8$ is the same one LIGO reads off a real binary, and the mesh reproduces it to the measured $-0.375$.

\subsection{The binding force is the self's, not gravity's}

A self must bind a body that still radiates. The mechanism that does it, binding by a well rather than by reflection, the lean's active vacuum mediating a deterministic Casimir attraction, and the reason a reflecting pin seals radiation in, is owned by the section on how a self holds itself together, with the one-tone matter-radiation framing in the matter-and-radiation section.

The capture recipe behind that binding is settled and measured. \emph{Capture} is an attractive interaction (a bound state to fall into) plus a bath (somewhere to shed the binding energy into). The bath is the open wake, where radiation that reaches the frontier is gone for good. The attraction is what was missing, and the search for it told an honest story.

\begin{result}[Capture needs both halves]
A pair launched apart in an attractive well coupled to a radiative bath inspirals and settles. The late amplitude falls to near zero, which is capture. Remove the bath (a reflecting end) and it oscillates forever without settling. Remove the attraction (a repulsive well) and it escapes. Both ingredients are necessary. Capture equals attraction plus bath.
\end{result}

The bare knit alone gives neither. Two charges on parallel non-colliding tracks evolve as the exact overlay of each alone, mismatch exactly zero, so the committed reversible knit exerts \emph{no force at a distance}. Its only interaction is contact, and contact scatters rather than binds. The attraction does not come free from the bare collide-then-stream law. It is supplied by the lean's active vacuum, the deterministic Casimir pressure that strengthens as bodies approach.

The point that matters for gravity is the boundary between the two. The instantaneous-Newtonian well above is an \emph{effective} attraction, slaved to mass and dissipative. The genuine reversible self-binding lives at the emergent scale, where binding and radiation act on different scales with no local contradiction. Gravity here is the effective field. The reversible binding is the self section's.

\subsection{Why curvature is essential}

A flat lattice, the Euclidean $\{3,4,3,3\}$ with the same $24$ sites, has no curvature. Hence no curvature-sourced gravity, no holographic boundary, no expansion. This is a clean control that isolates curvature as the source of the gravitational structure.

\begin{principle}[Curvature is the source]
Strip the negative curvature and keep everything else, the sites, the tone, the knit, and the gravitational structure vanishes. The flat control has no attraction from geometry, no boundary to write a hologram on, and no growing wake. Curvature is not decoration. It is the source.
\end{principle}

\subsection{The honest open frontier}

Two things are genuinely open, and we name them plainly.

\begin{enumerate}
\item \textbf{Curved-bulk gravity on the hyperbolic mesh.} Two attempts did not discriminate cleanly. The bulk is boundary-dominated, its shells grow exponentially, so bulk curvature disperses. It makes a stronger bath, not a clean binder. The clean $1/r$ Newton is a flat-cusp result, not a curved-bulk one. We do not overclaim curved-bulk gravity.
\item \textbf{Reversible discrete propagating gravity.} The effective field is dissipative. A stable reversible propagating version hits the discrete Klein-Gordon instability.
\end{enumerate}

What the mesh has is the \emph{effective theory plus quantitative relations}, the Einstein equation of state, the limits, the thermodynamics, the wave forms, all tied to mesh-measured inputs, the area law, the light cone, the flat cusp, the growth rate of the wake. The one thing not yet done by anyone from a micro-knit is full nonlinear numerical relativity simulated dock by dock. That is what "the mesh supports general relativity" honestly means.

\begin{principle}[The honest summary]
In the carrying image, gravity is how the fire's depth shapes the surface, the bulk drawing distant sparks together. The flat-cusp $1/r$ is solid. The curved-bulk pull is still smoke.
\end{principle}
